\section{Proofs and Further Details for Section~\ref{sec:acquired-growth}}\label{app:proofs-acquired}

Symbolic and numerical checks are in \texttt{code/verify\_part2\_acquirer.py}, section E.

\paragraph{Proof of \autoref{prop:acquired-score}.}
(a) $R^{\rm A}_X=X(g^{\rm org}+g^{\rm acq})+m-k^{\rm acq}g^{\rm acq}=R^{\rm org}_X+(X-k^{\rm acq})g^{\rm acq}$, whose derivative in $g^{\rm acq}$ is $X-k^{\rm acq}$. The score $Xg+m$ omits the term $-k^{\rm acq}g^{\rm acq}$. In \autoref{def:acquired-growth}, $\ln(S^{\rm cn}_t/S_{t-1})+\ln(1+S^{\rm ac}_t/S^{\rm cn}_t)=\ln(S_t/S_{t-1})$ because $S_t=S^{\rm cn}_t+S^{\rm ac}_t$, and dividing $S_t=S^{\rm cn}_t+S^{\rm ac}_t$ by $S_{t-1}$ gives the simple rates.

(b) In \autoref{ass:capital} each unit of growth costs $k$ of capital in the free-cash-flow margin; a unit of acquired growth costs $k^{\rm acq}$, so the margin after acquisitions is $m-k^{\rm acq}g^{\rm acq}$ and $R^{\rm A}_X=[Xg^{\rm org}+m]+(X-k^{\rm acq})g^{\rm acq}$. If acquisitions leave the growth technology of the businesses already owned unchanged, $g^{\rm org}$ and $m$ depend on spend $\sigma$ alone, as in \autoref{ass:growth}, so the problem separates: the first bracket is maximized at $R^*_X$ by \autoref{prop:score}, and the second is linear in $g^{\rm acq}\in[0,\bar g^{\rm acq}]$, maximized at $\bar g^{\rm acq}$ if $X>k^{\rm acq}$ and at 0 if $X<k^{\rm acq}$. Adding gives \eqref{eq:acquired-Rstar}.

(c) The operating margin $m_{op}=\gamma-a-\sigma$ of \autoref{ass:costs} and the free-cash-flow margin of \autoref{ass:capital} contain no financing cost. Financing an acquisition with debt changes $c_{\rm fin}$ and no term of $Xg+m_{op}$. A cash margin from which interest paid is deducted equals the margin before interest less the interest paid, net of any tax it saves, per dollar of revenue; under U.S.\ GAAP operating cash flow is such a margin, because interest paid is an operating cash outflow \citep{deloitte2026cashflow}.

(d) By definition, NRR is subscription revenue in a period from customers present a year earlier divided by subscription revenue a year earlier. Without new customers, subscription revenue therefore changes over the year by $(\mathrm{NRR}-1)S_{{\rm sub},t-1}$ and the rest of revenue by $\gamma_{\rm ns}(S_{t-1}-S_{{\rm sub},t-1})$, so revenue changes by the share
\[
a=\frac{S_{{\rm sub},t-1}}{S_{t-1}}(\mathrm{NRR}-1)+\Bigl(1-\frac{S_{{\rm sub},t-1}}{S_{t-1}}\Bigr)\gamma_{\rm ns}
\]
of the previous year's revenue, which is the base's contribution to annual organic growth. Converting the annual rate, the base grows at the continuously compounded rate $G_i=\ln(1+a)$ over the year, the zero-spend growth of \autoref{rem:assumptions}(b), under which $R_\infty=m_0+(X-k)G_i$. If $S_{\rm sub}=S$, or if $\gamma_{\rm ns}=\mathrm{NRR}-1$, then $a=\mathrm{NRR}-1$ and $G_i=\ln\mathrm{NRR}$. If $\mathrm{NRR}<1$ and $\gamma_{\rm ns}\in[\mathrm{NRR}-1,0]$, then $a$, a weighted average of $\mathrm{NRR}-1$ and $\gamma_{\rm ns}$, lies in $[\mathrm{NRR}-1,0]$, and since $\ln(1+\cdot)$ is increasing, $G_i\in[\ln\mathrm{NRR},0]$. Organic growth also includes the revenue of new customers of businesses already owned, which NRR does not measure.

(e) By \autoref{rem:assumptions}(b), \autoref{prop:budget}(a) holds with zero-spend growth $G_i$ and the right side
\[
(X-k)\bigl[\ln(s(T)/s_0)+(G-G_i)T\bigr].
\]
Applied to business $b$ in its own market from the date it is acquired, with initial share $s_b$, cap $\bar s_b$ and $s(T)<\bar s_b$, it gives the stated bound. If no two businesses share a market, the shares and caps are those of different markets, so the firm's organic budget is the sum of these bounds over the businesses it owns, and each acquisition adds a nonnegative term. If an acquired business sells in a market the firm already serves, its share adds to the firm's share of that market under the same cap, and the bound applies to the combined share.

(f) Revenue grows at $g=g^{\rm org}+g^{\rm acq}$, so holding $g$ requires $g^{\rm acq}=g-g^{\rm org}=g+d^{\rm org}$. Substituting into $R^{\rm A}_X=Xg+m-k^{\rm acq}g^{\rm acq}$ gives the expression, whose derivative in $d^{\rm org}$ at fixed $g$ is $-k^{\rm acq}$. If no new customers arrive, $g^{\rm org}=G_i$ by (d), so $g^{\rm acq}=g-G_i$, which exceeds $g$ when $G_i<0$, and $G_i=\ln\mathrm{NRR}$ when all revenue is subscription revenue. \qed

\paragraph{The inequalities of \autoref{rem:acquired-bounds}.}
$\mathrm{CRF}(r,T^{\rm life})\int_0^{T^{\rm life}}e^{-rt}\,dt=1$, so a level charge at that rate for $T^{\rm life}$ years repays a dollar at $r$; the factor falls with $T^{\rm life}$, exceeds $1/T^{\rm life}$ (since $1-e^{-x}<x$) and tends to $1/T^{\rm life}$ as $r\to0$. The user cost $(r+d)C^{\rm acq}_ve^{-d(t-v)}$ of an outlay whose value decays at $d$ is the return $r$ on its current value plus its depreciation, and $\int_v^\infty(r+d)e^{-d(t-v)}e^{-r(t-v)}\,dt=1$, so these charges also repay a dollar at $r$. Time-weighting replaces the value of capital at the end of year $t$ by its average over the year, $C^{\rm acq}_ve^{-d(t-1-v)}(1-e^{-d})/d$ for an outlay at $v\le t-1$ and $C^{\rm acq}_v(1-e^{-d(t-v)})/d$ for an outlay at $v\in(t-1,t]$, with $d\to0$ giving the undecayed weights $1$ and $t-v$. Every charge is nonnegative, so each charged score lies below the uncharged score, and it lies above the fully charged score if and only if the year's charge is at most $C^{\rm acq}_t$; for the level annuity on year-end capital that is $\mathrm{CRF}(r,T^{\rm life})\sum_{t-T^{\rm life}<v\le t}C^{\rm acq}_v\le C^{\rm acq}_t$. \qed

\subsection{Interest paid and events after the second-quarter report}\label{app:acquired-after}
This subsection accompanies Section~\ref{sec:acquired-growth}. The final prospectus reports interest paid, \acqIntPaidShTwentyThree, \acqIntPaidShTwentyFour{} and \acqIntPaidShTwentyFive{} of revenue in 2023, 2024 and 2025 and \acqIntPaidShQOne{} in the first quarter of 2026 \citep{bendingspoons2026prospectus}; in the first halves of 2025 and 2026, interest paid was \acqIntPaidRatioHOneFive{} and \acqIntPaidRatioHOneSix{} of interest expense \citep{bendingspoons2026sixk}. Adding back interest paid instead of interest expense moves the 2025 charged scores of Table~\ref{tab:acquired-score} to \acqChPaidLo{} to \acqChPaidHi.

The scores of Table~\ref{tab:acquired-score} are as of the second-quarter report of 13 August 2026. The company's reports on Form 6-K through 29 October 2026 record further acquisitions and financing \citep{bendingspoons2026sixk}. On 28 July it announced a EUR~500 million term loan guaranteed by SACE, Italy's export credit agency, which with EUR~495 million of additional term loan A financing and a EUR~490 million increase in revolving credit commitments brought the new and expanded facilities agreed since the start of the second quarter to EUR~1.49 billion, available for general corporate purposes and acquisitions. It bought Airtable, at an enterprise value of \$1.285 billion, with annual recurring revenue of about \$480 million as of June 2026 (agreed 4 August, completed 4 September), and Miro, at an enterprise value of \$1.355 billion, with annual recurring revenue of around \$600 million (agreed 10 September, completed 26 October), both for cash; at the Miro closing some of Miro's former shareholders invested about \$295 million of their proceeds in about 6.2 million new shares. These are enterprise values, net of the targets' cash, against annual recurring revenue rather than GAAP revenue, so they do not compare one for one with the prices per dollar of annual revenue below. It added \$1.25 billion and EUR~395 million to its term loans B due 2031 (priced 25 September, closed 2 October), upsized from targets of \$750 million and the equivalent of \$250 million, after Moody's raised its rating to Ba3 from B1 and S\&P to BB$-$ from B+; the proceeds were to finance Miro, repay the drawn part of its euro revolving credit facility and serve general corporate purposes. Berkshire Hathaway bought 10 million new shares at \$33.00 each, \$330 million, which with the \$600 million it invested at the listing took its stake to 4\% (announced 14 October, closed 29 October). And two of its subsidiaries priced \$1.25 billion of 9.625\% and EUR~500 million of 8.625\% senior secured notes due 2032, issued at par and upsized from about \$1.5 billion; the company does not issue or guarantee them, and the net proceeds are intended to fund future acquisitions (27 and 28 October, with settlement expected on or about 4 November). The new term loans and notes add to the debt behind the after-interest scores, which rest on the second quarter. Interest expense of the first half of 2026 at an annual rate, \$\acqRunRateIntHOne{} million, plus the coupons of the new notes, \$\acqNotesCoupon{} million, comes to \$\acqRunRateInt{} million a year, \acqRunRateShQ{} of second-quarter revenue at an annual rate; pricing the \$\acqUnpricedDebt{} million of new term loans, whose rates the reports do not state, at the notes' \acqCouponMarg\% gives \$\acqRunRateIntHi{} million, \acqRunRateShQHi. Adding the annual recurring revenue of Airtable and Miro, which is not revenue, to the denominator gives \acqRunRateShPf{} to \acqRunRateShPfHi. The add-ons were also to repay the drawn part of a revolving facility, so the upper case overstates. These ratios compare with the second quarter's \acqQTwoInt{} of revenue. The notes' stated use of proceeds shows continued reliance on acquisition capital, which is why a high charged score is evidence of access to targets and capital; it does not by itself show that the base is shrinking. We take no view on the company's valuation or on any shareholder's investment.
